\section{Torsion-free groups: trace rank and idempotents}
\label{sec:torsion-free}

For a torsion-free group, the support theorem leaves only the identity
conjugacy class. We first identify its coefficient with augmentation
rank, and then use the two possible scalar ranks to determine scalar
idempotents.

\begin{proof}[Proof of Corollary~\ref{cor:torsion-free}\textup{(i)}]
Let $e\in M_q(\C G)$ be idempotent and put
$x=[e(\C G)^q]\in K_0(\C G)$. Since $G$ is torsion-free,
Theorem~\ref{thm:bass} says that $\HS_G(x)$ is supported on $[1]_G$.
Therefore
\begin{align}
 \tau_{G,q}(e)
 &=\HS_G(x)([1]_G)\notag\\
 &=\sum_{C\in\con(G)}\HS_G(x)(C)\notag\\
 &=\epsilon_G\!\left(\sum_{i=1}^q e_{ii}\right)
 =\tr\bigl(\epsilon_{G,q}(e)\bigr)
 =\operatorname{rank}_{\C}\bigl(\epsilon_{G,q}(e)\bigr).
 \label{eq:torsion-free-rank}
\end{align}
The class sum is finite because the entries of $e$ have finite group
support. The last
equality holds because $\epsilon_{G,q}(e)$ is an idempotent complex
matrix: its image-kernel decomposition makes its trace equal to its
rank. Taking differences of projective classes gives
$\tau_{G,*}=\operatorname{rk}_{\epsilon}$. This map has range exactly
$\Z$, since its values are integers and $[\C G]$ has value $1$.
\end{proof}

For matrices the conclusion is numerical. For example,
$\operatorname{diag}(1,0)\in M_2(\C G)$ is an idempotent other than zero
and the identity, with augmentation rank one. The scalar case is special
because its augmentation rank can only be zero or one.

\begin{proof}[Proof of Corollary~\ref{cor:torsion-free}\textup{(ii)}]
First take $R=\C$ and let $e\in\C G$ be idempotent. Equation
\eqref{eq:torsion-free-rank} with $q=1$ gives
\[
 \tau_G(e)=\epsilon_G(e)\in\{0,1\}.
\]
We use the standard faithful-trace argument, as in
\cite[Theorem 2.1 and Remark 2.2]{BV1998}, spelling out a similarity
to a projection. Embed $\C G$ by left convolution in the group von
Neumann algebra $\mathcal N(G)$, the commutant of the right regular
representation on $\ell^2(G)$. Its canonical normalized trace is
\[
 \tau_G(T)=\langle T\delta_1,\delta_1\rangle .
\]
Its faithfulness can be seen directly: if
$T\geq0$ and $\tau_G(T)=0$, then $T^{1/2}\delta_1=0$. Since $T^{1/2}$
commutes with right translations and their translates of $\delta_1$
span a dense subspace, $T^{1/2}=0$, hence $T=0$.

Set
\[
 a=e^*e+(1-e)^*(1-e).
\]
For every $\xi\in\ell^2(G)$,
\[
 \langle a\xi,\xi\rangle
 =\|e\xi\|^2+\|(1-e)\xi\|^2
 \geq \tfrac12\|\xi\|^2.
\]
Thus $a$ is positive and invertible in $\mathcal N(G)$, and a direct
multiplication gives $ae=e^*a$. With $s=a^{1/2}$, the element
$p=ses^{-1}$ satisfies $p^2=p$ and
$p^*=s^{-1}e^*s=ses^{-1}=p$. Hence $p$ is a projection, and
traciality gives $\tau_G(p)=\tau_G(e)$. If this value is zero,
faithfulness applied to $p=p^*p$ gives $p=0$. If it is one,
faithfulness applied to $1-p$ gives $p=1$. Similarity then gives
$e=0$ or $e=1$.

Now let $R$ be any commutative unital domain of characteristic zero
and write an idempotent $f\in RG$ as
$f=\sum_{g\in S}r_g g$ with $S$ finite. Put
\[
 R_0=\Z[r_g:g\in S]\subseteq R,\qquad F=\operatorname{Frac}(R_0).
\]
The field $F$ is finitely generated over $\mathbb Q$, and therefore
embeds in $\C$. Indeed, send a finite transcendence basis of
$F/\mathbb Q$ to an algebraically independent tuple of complex
numbers; the resulting embedding of the rational function field
extends across the finite algebraic extension $F$ because $\C$ is
algebraically closed. The coefficientwise maps
$R_0G\hookrightarrow RG$ and $R_0G\hookrightarrow\C G$ are injective.
Thus $f^2=f$ already holds in $R_0G$, and its image in $\C G$ is an
idempotent. The complex case makes that image $0$ or $1$, and
injectivity gives the same conclusion for $f$. Only the finitely
generated coefficient field $F$ was embedded; no embedding of all
of $R$ into $\C$ is required.
\end{proof}

\begin{remark}[Algebraic and completed trace ranges]
Continue to assume that $G$ is torsion-free. Let
$j:\C G\hookrightarrow C_r^*(G)$ be the inclusion into the reduced
group $C^*$-algebra, and let
\[
 j_*:K_0(\C G)\longrightarrow K_0^{\mathrm{top}}(C_r^*(G))
\]
be the comparison map to topological $K$-theory. Write
$\tau_*^{\mathrm{top}}:K_0^{\mathrm{top}}(C_r^*(G))\to\R$ for the
canonical trace pairing. It agrees with $\tau_{G,*}$ on this image:
the same positive similarity used above, applied in a matrix algebra
over $C_r^*(G)$, produces a projection representing the image of a
matrix idempotent, and traciality preserves its trace. Consequently
\[
 \tau_*^{\mathrm{top}}\!\left(j_*K_0(\C G)\right)=\Z.
\]
This is the integer range on the algebraic image, not a statement
about every class of $K_0^{\mathrm{top}}(C_r^*(G))$. In particular,
the argument does not presume that the algebraic augmentation extends
to the reduced completion. The stronger completed idempotent question
is the Kadison--Kaplansky conjecture; see \cite[Section 1]{OW2023}.

The coefficient reduction above uses characteristic zero. An idempotent
$e\notin\{0,1\}$ would give nonzero zero divisors through $e(1-e)=0$,
so the zero-divisor conjecture would imply the idempotent conjecture.
The reverse implication is not established here.
\end{remark}
